\documentclass[a4paper,11pt,oneside]{memoir}

% Castellano
\usepackage[spanish]{babel}
\selectlanguage{spanish}
\usepackage[utf8]{inputenc}
\usepackage{placeins}

% Links
\usepackage[colorlinks]{hyperref}
\hypersetup{
	allcolors = {red}
}

% Ecuaciones
\usepackage{amsmath}

% Rutas de fichero / paquete
\newcommand{\ruta}[1]{{\sffamily #1}}

% Párrafos
\nonzeroparskip

% Imagenes
\usepackage{graphicx}
\newcommand{\imagen}[2]{
	\begin{figure}[!h]
		\centering
		\includegraphics[width=0.9\textwidth]{#1}
		\caption{#2}\label{fig:#1}
	\end{figure}
	\FloatBarrier
}

\newcommand{\imagenflotante}[2]{
	\begin{figure}%[!h]
		\centering
		\includegraphics[width=0.9\textwidth]{#1}
		\caption{#2}\label{fig:#1}
	\end{figure}
}

\graphicspath{ {./img/} }

% Capítulos
\chapterstyle{bianchi}
\newcommand{\capitulo}[2]{
	\setcounter{chapter}{#1}
	\setcounter{section}{0}
	\chapter*{#2}
	\addcontentsline{toc}{chapter}{#2}
	\markboth{#2}{#2}
}

% Apéndices
\renewcommand{\appendixname}{Apéndice}
\renewcommand*\cftappendixname{\appendixname}

\newcommand{\apendice}[1]{
	%\renewcommand{\thechapter}{A}
	\chapter{#1}
}

\renewcommand*\cftappendixname{\appendixname\ }

% Formato de portada
\makeatletter
\usepackage{xcolor}
\newcommand{\tutor}[1]{\def\@tutor{#1}}
\newcommand{\course}[1]{\def\@course{#1}}
\definecolor{cpardoBox}{HTML}{E6E6FF}
\def\maketitle{
  \null
  \thispagestyle{empty}
  % Cabecera ----------------
\noindent\includegraphics[width=\textwidth]{cabecera}\vspace{1cm}%
  \vfill
  % Título proyecto y escudo informática ----------------
  \colorbox{cpardoBox}{%
    \begin{minipage}{.8\textwidth}
      \vspace{.5cm}\Large
      \begin{center}
      \textbf{TFG del Grado en Ingeniería Informática}\vspace{.6cm}\\
      \textbf{\LARGE\@title{}}
      \end{center}
      \vspace{.2cm}
    \end{minipage}

  }%
  \hfill\begin{minipage}{.20\textwidth}
    \includegraphics[width=\textwidth]{escudoInfor}
  \end{minipage}
  \vfill
  % Datos de alumno, curso y tutores ------------------
  \begin{center}%
  {%
    \noindent\LARGE
    Presentado por \@author{}\\ 
    en Universidad de Burgos --- Curso \@course{}\\
    Tutor: \@tutor{}\\
  }%
  \end{center}%
  \null
  \cleardoublepage
  }
\makeatother


% Datos de portada
\title{PLQuiz}
\author{Roberto Izquierdo Amo}
\tutor{Dr. César Ignacio García Osorio}
\course{2013/2014}
\date{}

\begin{document}

\maketitle

\frontmatter

% Abstract en castellano
\renewcommand*\abstractname{Resumen}
\begin{abstract}
Desarrollo de una herramienta de escritorio que ayude a generar preguntas de test aleatorias (tipo quiz, cloze, de texto libre...) sobre problemas de algoritmos de análisis léxico.
El formato utilizado para generar las preguntas será importable a entornos virtuales de aprendizaje (Moodle), y permitirá su impresión en papel a través de la obtención de código \LaTeX .
\end{abstract}

\renewcommand*\abstractname{Descriptores}
\begin{abstract}
Generación de cuestionarios, análisis léxico, expresiones regulares, autómatas finitos.
\end{abstract}

\clearpage

% Abstract en inglés
\renewcommand*\abstractname{Abstract}
\begin{abstract}
Development of a desktop application that generates random test questions (quiz, cloze, free text...) about lexical analysis.
The format used to produce the questions should be compatible with virtual learning frameworks (such as Moodle), and allow for physical printing by generating \LaTeX\ code.
\end{abstract}

\renewcommand*\abstractname{Keywords}
\begin{abstract}
Questionnaire generation, lexical analysis, regular expresions, finite automata.
\end{abstract}

\clearpage

% Indices
\tableofcontents

\clearpage

\listoffigures

\clearpage

%\listoftables

%\clearpage

\mainmatter

\capitulo{1}{Objetivos del proyecto}

\section{Objetivos técnicos}
Los objetivos del proyecto, vistos desde el punto de vista del desarrollo técnico de la aplicación, pueden desglosarse en los siguientes puntos:
\begin{itemize}
	\item Desarrollo de una interfaz gráfica multiplataforma.
	\item Procesado y manipulación de expresiones regulares mediante el uso de un procesador de textos, mediante el uso de la herramienta \emph{JavaCC}.
	\item Resolución de problemas de aplicación del algoritmo de Aho-Sethi-Ullman a partir de las expresiones regulares procesadas.
	\item Resolución de problemas de aplicación del algoritmo de McNaughton-Yamada-Thompson a partir de las expresiones regulares procesadas.
	\item Agregación de los problemas en cuestionarios.
	\item Exportado de cuestionarios a dos formatos:
	\begin{itemize}
		\item Formato \LaTeX{}, generando cuestionarios listos para impresión, con intención de ser usados como examenes en papel.
		\item Formato \emph{XML} compatible con la plataforma de aprendizaje virtual Moodle.
	\end{itemize}
	\item Elaboración de un documento detallando los aspectos relevantes del desarrollo y funcionamiento de la aplicación.
\end{itemize}

Sin embargo los objetivos no han permanecido estáticos durante el periodo de desarrollo.
A los puntos iniciales se han añadido los siguientes:
\begin{itemize}
	\item Generación de problemas aleatorios, utilizando algoritmos de programación genética y métodos de búsqueda.
	\item Generación de problemas en bloque, produciendo cuestionarios completos con una variedad de problemas distintos.
	\item Exportado de imágenes en múltiples formatos, incluyendo exportado directo, exportado de programas \emph{Graphviz} y codificación en base 64.
	\item Exportado de cuestionarios a un formato adicional, HTML, permitiendo la vista previa de cuestionarios dentro de la interfaz gráfica.
	\item Implantado de un sistema de `log', facilitando el diagnóstico de problemas durante el desarrollo de la aplicación.
\end{itemize}

\section{Objetivos académicos}
Además de los objetivos relacionados directamente con el desarrollo del proyecto concreto, se consideran otros más generales, particularmente los detallados en la guía docente.
De estos destaco especialmente los siguientes:
\begin{itemize}
	\item Aplicación de los conocimientos teóricos y prácticos adquiridos en al asignatura al desarrollo de un proyecto de más envergadura y durante un periodo de tiempo considerable.
	\item Planteamiento y resolución de un problema real, con aplicaciones y usos concretos.
	\item Aprendizaje de nuevos temas relacionados con la titulación, trabajando con nuevos métodos y tecnologías y adaptando a ellos el modo de trabajo.
	\item Desarrollo de la capacidad de exposición y comunicación, especialmente durante la defensa, en el desarrollo de la memoria y en la preparación de la documentación de la aplicación.
\end{itemize}

\section{Objetivos personales}
En último lugar, la realización de este proyecto me ha dado la oportunidad de completar una serie de objetivos a nivel personal.
Destaco especialmente el aprendizaje de nuevas tecnologías y herramientas, que aunque ya me resultaban conocidas o familiares, no había tenido oportunidad de aplicar seriamente.

Entre las herramientas utilizadas en el proyecto, las utilizadas para el `logging', SLF4J y Logback, resultan mucho más útiles al ser aplicadas a bases de código grandes.
Su uso en el proyecto no solo ha facilitado el diagnóstico de errores y el control del flujo del programa, sino que me ha proporcionado conocimientos que probablemente pueda aplicar más adelante.

También destaco el uso de la versión 8 de Java, que apareció cuando el proyecto ya estaba avanzado.
Especialmente el añadido de los `stream', que permiten el manejo de datos de manera similar a la programación funcional, resultó de ayuda durante el desarrollo de ciertas partes de la aplicación.
Su uso en el proyecto me ha permitido ponerme al día de las novedades en Java, utilizándolas para resolver problemas reales.
\chapter*{Conceptos teóricos}
\capitulo{3}{Técnicas y herramientas}

\section{Técnicas}
Las técnicas empleadas en la aplicación y su desarrollo son las siguientes:

\subsection{SCRUM}
SCRUM\footnote{https://www.scrum.org/} es una metodología ágil que se define como
\begin{quote}
Un marco de trabajo por el cual las personas pueden acometer problemas complejos 
adaptativos, a la vez que entregar productos del máximo valor posible productiva y 
creativamente\cite{scrum}.
\end{quote}
Esta metodología se centra en crear procesos iterativos e incrementales, basandose en experiencias basadas para realizar estimaciones de rendimientos futuros, en base a las cuales tomar decisiones y controlar riesgos.

\subsection{UML}
UML\footnote{http://www.uml.org/} (\emph{Unified Modeling Language}) es un lenguaje de modelado utilizado en la estructura, comportamiento y arquitectura de aplicaciones.
El modelado de aplicaciones se considera una parte fundamental del diseño de las mismas, permitiendo examinar las funcionalidades desde un elevado nivel de abstracción.

\section{Herramientas}
Las herramientas utilizadas durante el desarrollo del proyecto son las siguientes:

\subsection{Java}
Java es un lenguaje de programación de propósito general \cite{jls8}, concurrente, basado en clases y orientado a objetos.
Se trata de un lenguaje fuerte y estáticamente tipado, y distingue claramente entre errores en tiempo de compilación y errores en tiempo de ejecución.

La compilación de un programa Java implica su traducción a representación `byte code', que es independiente de la plataforma.
Esta representación intermedia pasa a ejecutarse sobre la máquina virtual Java.
Por lo tanto, el código Java sigue la filosofía `write once, run anywhere', siendo compatible con cualquier arquitectura que disponga de una máquina virtual compatible.

Java es un lenguaje de relativo alto nivel, en el sentido de que los detalles de representación no son accesibles al programador.
Incluye una administración automática del almacenamiento de datos en forma del recolector de basura, diseñado para evitar los problemas relacionados con la liberación manual de memoria.

El lenguaje Java es, a fecha de 2014, uno de los lenguajes de programación más populares del mundo \cite{website:tiobe}.

\subsubsection{Java 8}
Java SE 8 es la edición más reciente del lenguaje de programación Java, y representa la mayor evolución del mismo en su historia.

Añade a la especificación un conjunto relativamente pequeño de nuevas características, mediante las cuales permite combinar modelos de programación orientada a objetos y funcional.
Estos cambios pretenden favorecer una serie de buenas prácticas --- inmutabilidad, ausencia de estado, composición --- a la vez que mantienen las características de Java --- legibilidad, simplicidad y universalidad.

Las librerías de la plataforma Java mantienen una evolución paralela a la del lenguaje.
Esto significa que el uso de las nuevas características --- expresiones lambda, referencias a métodos, interfaces funcionales --- es directamente compatible.
\cite{jls8}

\subsubsection{Eclipse}
El Eclipse Software Development Kit (Eclipse SDK)\footnote{http://eclipse.org/} es tanto un entorno de desarrollo integrado (IDE) para Java como una base para productos basados en Eclipse Platform.

Eclipse aparece como una herramienta propietaria de IBM que pretendía unificar los entornos de desarrollo ofrecidos a sus clientes y posibilitar la reutilización de componentes entre los mismos.
El proyecto de código abierto surge en 2001, con la fundación del Eclipse Consortium, que se encarga del desarrollo y administración del proyecto en la actualidad. \cite{website:eclipseFAQ}

Uno de los puntos fuertes de Eclipse es su capacidad para integrar una variedad de componentes mediante su sistema de plugins, lo cual permite trabajar con una variedad de lenguajes y herramientas.

\subsection{JUnit 4}
JUnit es un sistema software utilizado para realizar pruebas sobre código Java\footnote{http://junit.org/}, formando parte de la familia de herramientas de pruebas xUnit.
Es una de las librerías Java más utilizadas en proyectos de código abierto \cite{website:githubTOP}.

La versión 4 de JUnit extiende y simplifica la funcionalidad de anteriores versiones, haciendo un uso extensivo del sistema de anotaciones de Java.

\subsection{Apache Maven}
Maven es una herramienta de administración y construcción de proyectos usada principalmente con Java\footnote{http://maven.apache.org/}.
Similar en ciertos aspectos a Apache Ant, pero con un modelo de configuración más simple, y con el objetivo establecer una serie de `buenas prácticas' mediante sus configuraciones por defecto (`Convention over configuration').

Su funcionamiento se basa en la existencia de un fichero de configuración XML, el POM (Project Object Model).
Este fichero define la construcción, emisión documentación, empaquetado, pruebas, manejo de dependencias y múltiples otras tareas de manera centralizada \cite{mvnEx}.

Maven se estructura en torno a un núcleo central de tamaño mínimo, que puede ser extendido mediante \emph{plugins}.
Este núcleo es el encargado de descargar los \emph{plugins} que añaden la funcionalidad pedida desde sus repositorios remotos.

\subsection{SLF4J}
SLF4J (Simple Logging Facade for Java) sirve como interfaz para varias herramientas de logging\footnote{http://www.slf4j.org/}, permitiendo al usuario cambiar entre una y otra sin realizar modificaciones en el código.
Se establece como única dependencia obligatoria, encargándose por si misma de buscar una implementación compatible~\cite{website:slf4j}.

Esta herramienta permite trabajar con, por ejemplo, java.util.logging\footnote{http://docs.oracle.com/javase/8/docs/api/java/util/logging/package-summary.html}, logback\footnote{http://logback.qos.ch/} o log4j\footnote{http://logging.apache.org/log4j/1.2/}.

Es una de las librerías más utilizadas en proyectos de código abierto \cite{website:githubTOP}.

\subsection{Logback}
Logback es una herramienta de logging para Java que se establece como sucesor de Log4J.
Utiliza el API de SLF4J, por lo cuál puede ser libremente intercambiado por cualquier otro módulo de logging compatible.

La configuración de Logback se realiza mediante un pequeño fichero XML, en el cual se especifican los canales de salida, sean archivos o consola.
Siguiendo la estructura de Apache Maven, Logback permite establecer ficheros de configuración distintos para las etapas de desarrollo y para el producto final.
Logback refresca la configuración en tiempo de ejecución, permitiendo realizar cambios `en caliente' \cite{website:logback}.

Es la herramienta de logging más utilizada en proyectos de código abierto \cite{website:githubTOP}.

\subsection{JavaCC}
JavaCC\footnote{https://javacc.java.net/} (Java Compiler Compiler) es un generador de analizadores sintácticos Java.
Funciona mediante la especificación de una gramática en un formato propio.
A partir de este fichero, la herramienta genera un programa Java capaz de procesar texto y reconocer coincidencias con la gramática.

Además de generador en sí, JavaCC proporciona una serie de herramientas relacionadas, como un constructor de árboles (JJTree) \cite{website:javacc}.

\subsection{JGraphX}
JGraphX\footnote{https://github.com/jgraph/jgraphx} es una biblioteca de visualización de gráficos para Java de código abierto, derivada de la implementación comercial en JavaScript (mxGraph).
En las versiones anteriores la herramienta se conocía como JGraph \cite{website:jgraphx}.

Implementa una serie de funcionalidades básicas para trabajar con grafos, incluyendo visualización e interacción con los mismos.
De cara a aplicaciones Java, es compatible con Swing y permite exportar las imágenes generadas.

\subsection{Graphviz}
Graphviz es una herramienta de visualización de gráficos de código abierto\footnote{http://www.graphviz.org/}.
Proporciona una manera de representar información estructurada como diagramas, grafos y redes.
Sus aplicaciones se encuentran en arquitectura de redes, bioinformática, ingenieria del software, diseño de bases de datos y web, machine learning y diseño de interfaces gráficas \cite{website:graphviz}.

La herramienta toma descripciones de gráficos en texto plano y los transforma en diagramas, con varios formatos disponibles.
El formato `dot', en concreto, permite dibujar grafos dirigidos.

\subsection{Git}
Git es un sistema de control de versiones distribuido\footnote{http://git-scm.com/}, diseñado para trabajar con proyectos de cualquier tamaño.
Es un proyecto libre y de código abierto, licenciado bajo la GPL\footnote{http://www.gnu.org/copyleft/gpl.html}.

Una de las características distintivas de Git es su modelo de ramas.
Git permite mantener múltiples ramas locales independientes, cuya creación, modificación y combinado resulta poco costoso.
Esto permite aislar tareas, trabajando e incorporando cada una de manera totalmente separada.

Git es un sistema distribuido.
Esto quiere decir que el repositorio local mantiene una copia completa del repositorio remoto.
Esto permite una variedad de métodos de trabajo, y elimina puntos débiles del sistema \cite{website:git}.

\subsubsection{Github}
Github\footnote{https://github.com/} proporciona un servicio de almacenamiento de repositorios remotos y un entorno de colaboración para desarrolladores.
Actualmente cuenta con más de seis millones de usuarios y 13.5 millones de repositorios \cite{website:github}.

Github proporciona herramientas de seguimiento de proyectos, incluyendo una wiki y un issue tracker por repositorio.
También es compatible con otras aplicaciones web, como por ejemplo Pivotal Tracker, que facilitan su integración en el proceso de desarrollo.

\subsection{Pivotal tracker}
Pivotal Tracker\footnote{https://www.pivotaltracker.com} es una herramienta web para el manejo de proyectos ágiles, permitiendo la colaboración del equipo y la administración de tareas.
Su API permite la integración con otras herramientas de desarrollo, como por ejemplo Github.

Pivotal Tracker cuenta con una serie de herramientas de generación de informe, con las que podemos ver las tareas finalizadas en cada iteración, gráficos de \emph{burn-down}, etc.

\subsection{Moodle}
Moodle (\emph{Modular Object-Oriented Dynamic Learning Environment}) es una plataforma de enseñanza virtual (\emph{e-learning}) desarrollada como software libre.
Disponible en mas de 100 idiomas, Moodle habilita mas de 60 mil páginas de aprendizaje, con más de siete millones de cursos, 70 millones de alumnos y un millón de profesores \cite{website:moodle}.

Una de las características de Moodle es su capacidad para importar cuestionarios a partir de varios formatos, incluyendo texto plano, el formato propietario Gift y a partir de documentos XML.

\subsection{\LaTeX{}}
\LaTeX{} es un lenguaje de marcadores para la preparación de documentos comúnmente usado en publicaciones técnicas o científicas.

Es un procesador de texto «\emph{What you see is not what you get}», lo cual significa que lo que vemos durante la edición no es el 
documento final, sino las instrucciones para generarlo.
Esto permite separar casi completamente el contenido del documento de su formato.

\subsection{XML}
\emph{XML} (\emph{Extensive Markup Language}) es un lenguaje de etiquetas que se utiliza para crear documentos estructurados \cite{website:xml}, compuestos de entidades que contienen en su interior datos u otras entidades.
El estandar fue producido y es desarrollado por el \emph{World Wide Web Consortium}\footnote{http://www.w3.org/}.

Podemos verificar que un documento \emph{XML} tiene el formato correcto validandolo contra su esquema o \emph{DTD} (\emph{Document Type Definition}).
También podemos verificar que un documento está `bien formado', es decir, que cumple una serie de reglas gramaticales mínimas.

\subsection{ObjectAid UML Explorer}
ObjectAid UML Explorer\footnote{http://www.objectaid.com/} es un \emph{plugin} de visualización de código para eclipse.
Permite mostrar el código fuente de un proyecto Java en forma de diagramas UML, reflejando el estado y las relaciones en el código, y actualizandose a medida que el código cambia.
\capitulo{4}{Aspectos relevantes del desarrollo del proyecto}

\section{Desarrollo de prototipo}
En el intervalo de tiempo situado entre la selección del proyecto y el comienzo formal de la primera iteración de desarrollo, se dedicaron recursos al desarrollo de un prototipo de la aplicación.
Los objetivos de esta fase eran la construcción de una interfaz gráfica que aproximase el diseño final (aunque no todas las funciones estuvieran implementadas), y el desarrollo de una funcionalidad parcial de la aplicación.
En este caso la funcionalidad parcial era la resolución del primer tipo de cuestiones.

\imagen{prototipo_interfaz}{Interfaz gráfica del prototipo de la aplicación}

El prototipo está limitado a trabajar con un único problema a la vez, las opciones de generación no son funcionales y el botón `Generar' produce siempre la misma expresión.
A pesar de estas limitaciones, el prototipo aproxima con bastante precisión la interfaz final, y forma la base de resolución de problemas de manera eficaz.
Es capaz, en concreto, de procesar una expresion regular cualquiera proporcionada como texto, construir a partir de ella un árbol sintáctico, aplicarle el algoritmo de Aho-Sethi-Ullman y mostrar los resultados obtenidos.

\section{Metodología ágil - Scrum}
El desarrollo del proyecto ha sido enmarcado dentro de una metodología ágil, Scrum\footnote{\url{https://www.scrum.org/}}.
Sin embargo, y debido a las diferentes condiciones entre equipo de desarrollo y las condiciones entre este proyecto y uno más habitual, se han realizado ciertas modificaciones.

En cuanto a roles, volcamos los tres roles que define Scrum («\emph{Product Owner}», «\emph{Development Team}» y «\emph{Scrum Master}») sobre el desarrollador, dado que no tenemos un equipo como tal.
El rol de «\emph{Product Owner}», especialmente la definición de tareas para el \emph{backlog}, se realizará conjuntamente con el tutor.

Se definen \emph{sprints} de dos semanas, entre los cuales se realiza una reunión con el tutor del proyecto, y al final de cada cuál se completa un entregable.
Esta reunión cumple las funciones de «\emph{Sprint Planning Meeting}» y «\emph{Sprint Review}», y en ella se revisan los resultados del \emph{sprint} que termina y se definen los objetivos del que comienza.
No se realizan reuniones diarias («\emph{Daily Scrum}»), resultando innecesarias al contar con un equipo de desarrollo de una persona.

El seguimiento de los sprint se realiza mediante el uso de una herramienta online, Pivotal Tracker\footnote{\url{https://www.pivotaltracker.com/s/projects/1026880}}, de manera que resulta accesible a tutor y desarrollador.
La integración de esta herramienta con el sistema de control de versiones utilizado, Github\footnote{\url{https://github.com/RobertoIA/PLQuiz}}, permite la terminación de tareas directamente desde el \emph{commit} que las finaliza.
Se consideran validadas las tareas cuando se completa su documentación (en forma de \emph{javadoc}), se escriben sus pruebas y estas se ejecutan correctamente.

En el apéndice C describimos en detalle el desarrollo de cada \emph{sprint}, asi como su duración, objetivos y resultados.

\section{Estructura de paquetes}
La filosofía seguida a la hora de diseñar la estructura de la aplicación se centra en los conceptos de modularidad y reusabilidad.
Cada paquete pretende ser independiente, y las dependencias externas se reducen lo más posible.
Esto se traduce en un uso frecuente del patrón `fachada'.

Cada paquete principal tiene una estructura de dos niveles: el primer nivel contiene las clases `fachada' --- visibles al resto de la aplicación ---, mientras que el segundo (\emph{.datos}) contiene la lógica y estructuras de datos internas --- utilizadas solo por el paquete `padre' ---.

Podemos dividir los paquetes principales de la aplicación en tres clases:
\begin{itemize}
	\item Paquete de interfaz gráfica (\emph{es.ubu.inf.tfg.ui}).
	Contiene las clases que definen la apariencia y comportamiento de la interfaz de usuario.
	Es, por pura necesidad, el paquete con más dependencias externas, dependiendo tanto del paquete que define la representación de documentos como de los paquetes que definen cada clase de problemas.
	\item Paquete de construcción de documentos (\emph{es.ubu.inf.tfg.doc}).
	Dependen, por necesidad, de los paquetes que definen las clases de problemas que van a representar.
	\item Paquetes de problemas (\emph{es.ubu.inf.tfg.regex.thompson} y \emph{es.ubu.inf.tfg.regex.asu}).
	Contienen las clases encargadas de representar las estructuras internas del problema, sus soluciones, y aquellas necesarias para generarlos.
	Dependen únicamente del paquete de representación de expresiones regulares, y son independientes el uno del otro.
	\item Paquete de representación y procesamiento de expresiones regulares (\emph{es.ubu.inf.tfg.regex.parser} y \emph{es.ubu.inf.tfg.regex.datos}).
	Contienen las clases necesarias para procesar expresiones regulares a partir de cadenas de caracteres, representarlas y operar sobre ellas.
	Son la base sobre la cual trabajan los problemas.
\end{itemize}

\imagen{diagrama_paquetes}{Diagrama de paquetes y sus dependencias}

Como podemos ver, los paquetes encargados del procesamiento de expresiones regulares son la base del problema.
Junto con los paquetes de problemas forman el núcleo de la aplicación, y pueden ser utilizados de manera independiente, al carecer de dependencias externas.

Sobre estos paquetes centrales se aposentan los paquetes de representación de documentos, que toman los datos de los documentos y les dan un formato adecuado.
En un nivel superior esta la interfaz de usuario, que toma el conjunto de datos obtenido del resto de paquetes y lo presenta al usuario.

Esta estructura en niveles significa que podemos tomar cualquier paquete, y utilizarlo de manera completamente independiente, siempre que dispongamos de aquellos que se encuentran en niveles inferiores.

\section{Java 8 - \emph{expresiones lambda} y \emph{streams}}
Dos de las novedades introducidas con la versión 8 de Java son los \emph{streams}\footnote{\url{http://docs.oracle.com/javase/8/docs/api/java/util/stream/package-summary.html}} y las \emph{expresiones lambda}\footnote{\url{http://docs.oracle.com/javase/tutorial/java/javaOO/lambdaexpressions.html}}.

Los \emph{streams} permiten manipular colecciones, aplicándoles operaciones típicas de programación funcional, como pueden ser \emph{map} o \emph{reduce}.
Una de las ventajas de estas operaciones son que permiten la aplicación de operaciones a cada elemento de una colección sin iterar explicitamente, o la obtención de subconjuntos de la misma que cumplan una condición cualquiera.

La aplicación de operaciones se realiza mediante \emph{expresiones lambda}.
Las \emph{expresiones lambda} permiten definir una clase de un solo método de manera compacta, encapsulado la operación.

Dado que la salida de la nueva versión de Java ocurrió cuando el desarrollo de la aplicación ya estaba avanzado, la utilización de las nuevas herramientas se da de manera ocasional, y tiene un impacto reducido.

El principal uso de los \emph{stream} se encuentra en las clases encargadas de la generación de expresiones regulares.
Por ejemplo, las operaciones sobre `poblaciones' de expresiones regulares dentro de un algoritmo genético resultan simples si disponemos de medios de `filtrar' subconjuntos dentro de la colección.
Usando la función de evaluación del algoritmo genético como argumento para la operación de filtrado sobre el \emph{stream} permite obtener el conjunto de la población que queremos preservar, mutar o cruzar de manera simple y sin recurrir a la iteración explicita sobre la colección.
\capitulo{5}{Trabajos relacionados}
\capitulo{6}{Líneas de trabajo futuras}

\section{Algoritmos de generación de árboles de expresión regular}
El sistema de generación de expresiones regulares actual implementa el método \emph{full} descrito por John R. Koza \cite{koza92}.
Dicho método genera árboles en los que todas las hojas tienen la misma altura.

Este método permite generar expresiones de forma rápida y resulta relativamente simple de implementar, dado que los árboles no varían en profundidad.

Como desventaja, el método \emph{full} genera un rango de expresiones limitado.
Esta característica supone un problema a la hora de utilizar las expresiones para generar problemas.

Si agrupamos los problemas según los valores de las características que encontramos para generarlos, podemos ver que la distribución de los mismos no resulta uniforme, sino que se generan muchos problemas con ciertos valores y muy pocos o ninguno con otros.
Esto resulta problemático si presentamos una interfaz de usuario que permita introducir valores arbitrarios.

En la literatura encontramos definidos otros dos posibles algoritmos, \emph{grow} y \emph{half-and-half} \cite{koza92}.
La implementación de uno de estos algoritmos (o de algún otro equivalente) supondría una mejora del sistema de generación y permitiría comprobar como de equilibrada está la distribución de resultados.

Las mejoras en el algoritmo de generación de árboles afectarían también a los algoritmos de búsqueda, permitiendo tal vez utilizar métodos más rápidos o eficientes, y generar problemas más complejos.

\section{Mejoras a la interfaz gráfica}
\subsection{Opción de guardado y carga de cuestionarios}
En la versión actual de la aplicación, los cuestionarios pueden exportarse pero no importarse.
Esto quiere decir que una vez se cierra el programa no puede continuar trabajándose sobre las preguntas que teníamos, a menos que se cree un cuestionario nuevo y se añadan las preguntas a mano.

Se proponen dos posibles implementaciones de un sistema de almacenamiento y carga:
\begin{itemize}
	\item El método más evidente es el almacenamiento de los cuestionarios como ficheros separados, ya sea en texto plano, \emph{XML} o cualquier otro formato.
	Lo único que necesitamos conocer para reconstruir un problema es la expresión regular que lo define, y el tipo y sub-tipo de problema con el que se resuelve.
	Por lo tanto la cantidad de información a almacenar es muy reducida, y el sistema que la almacene y recupere sería simple de implementar.
	\item Un método más complejo, pero que no requiere guardar ficheros separados, consistiría en la lectura de un fichero exportado y su comparación con la plantilla de la que se generó.
	De esta manera podemos localizar la expresión y extraerla.
	Una vez obtenidas las expresiones regulares e identificado el tipo de problema al que pertenece cada una, la reconstrucción del cuestionario es trivial.
\end{itemize}

\subsection{Opción de deshacer cambios}
La aplicación no permite `volver atrás' si eliminamos parte del trabajo realizado, como por ejemplo eliminando un problema del cuestionario, o pulsando el botón de `Generar' en un problema que queríamos conservar.
Una opción de `deshacer' evitaría problemas al usuario y aumentaría la usabilidad del programa.

\subsection{Opción de reordenar cuestiones}
Las cuestiones se muestran en el orden en que se añadieron, sin permitir cambios.
Si se pretende añadir una cuestión entre varias existentes primero deben eliminarse las que se encuentran por debajo de la nueva posición.

La implementación del cambio de orden es sencilla, ya que el documento que la aplicación utiliza internamente ya almacena los problemas de manera ordenada.
Es necesario añadir un sistema que permita especificar un nuevo orden en el documento o una nueva posición para un problema dado.
Una posibilidad sería asociar cada problema con un número que indique su posición.

Se proponen varias implementaciones de cara a la interfaz gráfica:
\begin{itemize}
	\item La manera más sencilla de implementar la ordenación sería añadir dos botones al `frame' que representa el problema, uno para mover el problema hacia arriba y otro para moverlo hacia abajo.
	Es una implementación con usabilidad limitada, ya que tenemos que pulsar el botón tantas veces como posiciones queramos mover el problema.
	\item Una implementación sencilla y más usable sería el añadir un cuadro de texto o lista desplegable al `frame' del problema, que nos permita seleccionar una nueva posición para el problema.
	\item La implementación más compleja, pero más adecuada desde el punto de vista de la usabilidad, es la implementación de un sistema de `drag and drop'.
	Es decir, que la aplicación nos permita arrastrar un problema hasta su nueva posición.
	El problema de esta idea es que \emph{Swing} no soporta estas operaciones, y añadirlas supondría un esfuerzo de desarrollo considerable.
\end{itemize}

\section{Mejora del sistema de plantillas}
\subsection{Personalización de plantillas}
El sistema de plantillas de la aplicación permite la modificación de las mismas para conseguir documentos personalizados.
Sin embargo, la personalización de las plantillas requiere modificación y empaquetado del proyecto, o modificación directa de los ficheros del \emph{.jar}.

Una posible mejora sería la opción de proveer plantillas personalizadas.
Estas plantillas se seleccionarían al realizar la exportación del cuestionario, o desde un menú de opciones en la propia aplicación.
Las plantillas por defecto seguirían estando disponibles, en caso de que no se haya creado una propia y como ejemplo.

Con esta mejora no solo permitimos personalizar el estilo, sino también el contenido.
El usuario puede añadir o eliminar etiquetas, haciendo que el problema disponga de más o menos información.

\subsection{Integración de un motor de plantillas}
La herramienta de plantillas que la aplicación utiliza es propia, simple, y se basa en sustitución mediante expresiones regulares.
Una posible mejora es la introducción de un motor de plantillas externo en forma de librería, que aumente las opciones a la hora de trabajar con las mismas, o que las simplifique.

\section{Más tipos de cuestiones}
Una de las mayores limitaciones de la aplicación es el reducido número de tipos de problemas de los que dispone, y a partir de los cuales se generan los cuestionarios.
Una de las mejoras más obvias es una ampliación en el tipo de cuestiones disponibles.

\subsection{Integración con BURGRAM}
La aplicación BURGRAM\footnote{Carlos Gómez Palacios, Enero 2008, Universidad de Burgos}, desarrollada como proyecto de fín de carrera, permite la resolución de problemas asociados a gramáticas tipos \emph{LL}, \emph{LR}, \emph{SLR} y \emph{LALR}.
Aunque estos tipos de problemas son bastante diferentes a los que la aplicación resuelve actualmente, se enmarcan también dentro del área de procesamiento del lenguaje.
Los consideramos, por lo tanto, un añadido interesante a la aplicación.

La aplicación BURGRAM trata la resolución de los problemas paso a paso y a partir de una gramática dada.
El procesamiento y resolución del problemas están ya resueltos, por lo que su integración consistiría en adaptar la interfaz, añadir los distintos tipos al modelo de documentos y preparar las plantillas.
Idealmente la integración trataría a BURGRAM como una librería externa, sin incluir directamente su código dentro del proyecto.

\subsection{Compatibilidad con ficheros\emph{GIFT}}
El formato \emph{GIFT} es el usado por la plataforma de aprendizaje virtual Moodle para el almacenamiento de cuestiones en formato de texto plano.
Los tipos de preguntas disponibles son:
\begin{itemize}
	\item Selección múltiple.
	\item Verdadero o falso.
	\item Respuesta corta.
	\item Rellenar huecos.
	\item Cuestiones numéricas.
\end{itemize}

La compatibilidad con este formato nos permitiría importar o incluir con la aplicación un conjunto de preguntas de tipo teórico, que podrían insertarse en los cuestionarios.
Esto permitiría la creación de cuestionarios y examenes más completos, con preguntas de test, teóricas y resolución de problemas.

\appendix

\apendice{Generación de expresiones regulares}

Trabajar con una expresión regular dentro de la aplicación implica traducirla a una estructura de datos adecuada.
Contando con que la expresión es introducida como una cadena de caracteres, utilizaremos un parser para convertirla en un árbol binario.
Los algoritmos de generación de árboles son un tema muy estudiado dentro de la programación genética, lo cuál nos permite generar la expresión directamente en forma de árbol implementando un algoritmo ya conocido.
\\
El algoritmo concreto implementado es el del método "full" \cite{koza92}.
Este método involucra la creación de árboles cuya profundidad viene definida por la longitud del camino entre un extremo cualquiera y la raíz.
Esto quiere decir que generaremos árboles binarios en los que cada cada nodo no hoja tiene 1 o 2 hijos, y en el que todas las hojas se encuentran a la misma profundidad.
\\
El algoritmo funciona restringiendo la selección de etiquetas para los nodos generados en función de su profundidad.
Un nodo símbolo, por ejemplo, solo estará permitido en la profundidad máxima, mientras que un nodo cerradura estará permitido en cualquier profundidad excepto la máxima.
Tomaremos consideraciones adicionales a la hora de generar los hijos de un nodo, ya que el número de los mismos variará en función de la operación que contenga el padre.
Un nodo cerradura, por ejemplo, tiene un solo hijo, mientras que un nodo concatenación tendrá dos.

\section{Generación de problemas}

A la hora de generar problemas tomamos tres criterios:
\begin{itemize}
	\item El número de símbolos presentes en la expresión regular (sin contar el vacío).
	\item El número de estados en su función de transición.
	\item La presencia o ausencia del símbolo vacío.
\end{itemize}
Tanto el número de símbolos como la presencia del símbolo vacío dependen directamente de la expresión regular a partir de la cual construimos el problema.
El número de estados de la función de transición será el mismo para cualquier problema de un mismo tipo que comparta la misma expresión regular.
Podemos decir, por lo tanto, que las características que buscamos en un problema dependen exclusivamente de la expresión regular con la cuál lo construimos.
\\
La generación de problemas por lo tanto consistirá en generar expresiones regulares cuyo problema asociado cumpla las características pedidas.
Siendo el objetivo la generación de problemas con unas características dadas, deben determinarse los algoritmos de búsqueda a utilizar y los parámetros óptimos para los mismos.

\section{Coste de generación}

Dado que los árboles de expresión son árboles binarios, sabemos que el número máximo de nodos está acotado superiormente por la n
\[
expresi\acute{o}n\acute{u}mero \; de \; nodos \; \leq \; 2^{\;profundidad} - 1
\]
La generación de una expresión regular crea los nodos del árbol de la expresión uno a uno, por lo tanto el coste de generar una expresión crecerá con el número de nodos de la misma.
\\
Para comprobar estos datos de manera experimental obtendremos los tiempos medios de generación de 100.000 expresiones regulares para cada profundidad entre 0 y 8, y comparamos su crecimiento con el número máximo de nodos asociado a cada profundidad.

\imagen{profundidad}{Tiempo de generación de expresiones y número de nodos}

Podemos ver como el tiempo de generación de una expresión regular depende del número de nodos que contenta y, por lo tanto, de la profundidad.
Por lo tanto será preferible generar expresiones lo menos profundas posible, en concreto limitándonos a una profundidad máxima de 6.
\\
Destacamos que los valores del tiempo dependen de la máquina en que se realizan las pruebas.
Este análisis pretende analizar el crecimiento de la función tiempo, no sus valores exactos.

\section{Regiones aceptables}

Consideramos expresiones 'útiles' aquellas que contienen entre 2 y 6 símbolos distintos (sin incluir el vacío), y entre 3 y 15 estados distintos al generar a partir de la expresión un problema de tipo Aho-Sethi-Ullman o de construcción de subconjuntos.
Estos rangos corresponden con los problemas que consideramos apropiados para resolver de manera manual, y con lo permitido por la interfaz.
El generador es capaz de trabajar con rangos arbitrarios, pero los problemas que se forman con expresiones fuera de estos rangos son o demasiado complejos para resolver de manera manual, o completamente triviales.
\\
El conocimiento de que profundidades de árbol corresponden con más expresiones generadas dentro de los rangos deseados nos permite establecer los límites óptimos para utilizar con los generadores de problemas.
\\
Este análisis pretende determinar cuantas de las expresiones regulares generadas aleatoriamente encajan dentro de los rangos deseados, y en que medida depende de la profundidad del árbol sintáctico.
La intuición nos dice que los árboles más profundos corresponderán con las expresiones más complejas (más símbolos y/o estados).

\subsection{Aho-Sethi-Ullman}

Para realizar las pruebas generaremos 100.000 expresiones regulares conteniendo el símbolo vacío, y otras tantas sin contenerlo, para cada profundidad de árbol de entre 1 y 6.
Con cada una de las expresiones construiremos un problema de tipo Aho-Sethi-Ullman y comprobaremos cuantas de las expresiones generadas entran dentro de los rangos de símbolos y estados que consideramos aceptables.

\imagen{prof-ASU-NV}{Aho-Sethi-Ullman (no vacío) aceptables según profundidad}
\imagen{prof-ASU-V}{Aho-Sethi-Ullman (vacío) aceptables según profundidad}

Los resultados muestran que los árboles de profundidad 4 generan el mayor número de expresiones en la zona deseada.
También podemos ver claramente que las expresiones con árbol sintáctico de profundidad entre 3 y 6 generan una mayoría de expresiones regulares aceptables, tanto si incluimos el símbolo vacío como si no.
\\
Por lo tanto el método de búsqueda para problemas de Aho-Sethi-Ullman utilizará expresiones de profundidad entre 3 y 6.

\subsection{Construcción de subconjuntos}

Para el análisis de problemas de construcción de subconjuntos repetimos el mismo proceso que en el apartado anterior, generando expresiones regulares para cada combinación de tipo y profundidad, y construyendo problemas de construcción de subconjuntos con ellas.

\imagen{prof-CS-NV}{Construcción de subconjuntos (no vacío) aceptables según profundidad}
\imagen{prof-CS-V}{Construcción de subconjuntos (vacío) aceptables según profundidad}

Según los resultados podemos ver que las expresiones dentro del rango tienen mayoritariamente profundidades de entre 1 y 5, y que las expresiones de profundidad 6 están completamente fuera del rango.
\\
Por lo tanto el método de búsqueda para problemas de construcción de subconjuntos utilizará expresiones de profundidad entre 1 y 5.

\section{Distribución de resultados aceptables}

Partiendo de los resultados que encontramos dentro de la región aceptable, es importante identificar como se distribuyen esos resultados.
Asumimos que no todas las combinaciones de símbolo y estado van a aparecer con la misma frecuencia dentro de la región aceptable.
\\
Dado que la frecuencia de aparición de expresiones en la región aceptable varía según la profundidad, examinaremos cada profundidad utilizada por separado.
De esta manera podremos determinar si ciertas combinaciones de símbolo y estado son más probables en ciertas profundidades.
\\
Los datos experimentales utilizados son los mismos que en los apartados anteriores.

\subsection{Aho-Sethi-Ullman}

Agrupando los datos según profundidades podemos ver nuevamente las agrupaciones según profundidad.
Más interesante resulta que las expresiones generadas se agrupan, independientemente de la profundidad, alrededor de los mismos puntos.
\\
Es interesante notar que, aunque las agrupaciones no depende de la profundidad, las profundidades mayores parecen agruparse mas hacia la izquierda.
Esto implicaría que las profundidades mayores generan problemas con menos estados y símbolos o, más probablemente, que nuestro rango aceptable es demasiado restringido para ver la perspectiva completa.

\imagen{dist-ASU-NV}{Distribución de Aho-Sethi-Ullman (no vacío)}
\imagen{dist-ASU-V}{Distribución de Aho-Sethi-Ullman (vacío)}

Vemos que el comportamiento es similar entre expresiones que contienen el elemento vacío y aquellas que no.
\\
Esto indica, de manera bastante clara, que ciertas combinaciones de número de símbolos y estados son mucho más probables.
O, por otra parte, que en las profundidades con las que estamos trabajando tienden a generarse expresiones con unas características especificas.
\\
Por ejemplo, es posible que la combinación de 3 símbolos y 8 a 14 estados se dé con alta probabilidad en profundidades de 6 o mayores, o que necesite expresiones regulares con árboles que tengan hojas a distintas profundidades.

\subsection{Construcción de subconjuntos}

En los problemas de construcción vemos los mismos problemas que en los de Aho-Sethi-Ullman.
Las expresiones regulares se agrupan en torno a unas ciertas combinaciones de número de símbolos y estados, aunque en este caso las combinaciones son diferentes que en el apartado anterior.

\imagen{dist-CS-NV}{Distribución de construcción de subconjuntos (no vacío)}
\imagen{dist-CS-V}{Distribución de construcción de subconjuntos (vacío)}

Es interesante remarcar que las agrupaciones parecen ampliarse hacia la derecha.
Es decir, cuanto más elevado el número de símbolos, más variedad de estados generan los problemas.

\section{Algoritmos de búsqueda}

Una vez definidos los parámetros con los cuales vamos a generar las expresiones, queda pendiente encontrar un algoritmo de búsqueda que nos permita encontrar los problemas dados.
\\
Es probable que dos problemas generados a partir de expresiones regulares tengan características similares, pero no tenemos ninguna garantía al respecto.
Es decir, es posible que dos expresiones vecinas tengan características bastante diferentes.
Descartamos, por lo tanto, algoritmos de búsqueda como el recocido simulado, que se desplazan por el espacio de búsqueda de vecino a vecino.

\subsection{Búsqueda aleatoria}

El algoritmo de búsqueda aleatoria consiste en generar soluciones completamente al azar hasta encontrar la correcta.
Lo tomamos como posibilidad dado que resulta muy sencillo de implementar, y de que resultará eficiente en el caso de búsquedas sencillas.

% ASU, CS

\subsection{Algoritmo genético}

% ASU, CS, Reproducción o mutación

\section{Conclusiones}

% puede ser que la forma de los árboles de expresiones (nodos hoja con misma profundidad) favorezcan ciertas combinaciones de simbolos / estados y desfavorezcan otras.
\apendice{Manuales}

\section{Manual de usuario}

\subsection{Requerimientos}
La aplicación cuenta con los siguientes requerimientos:
\begin{itemize}
	\item Java 8\footnote{\url{http://www.oracle.com/technetwork/java/javase/downloads/jre8-downloads-2133155.html}}
	\item Windows
	\begin{itemize}
		\item Windows 8 (Desktop) / Windows 7 / Windows Vista SP2 / Windows Server 2008 / Windows Server 2012 (64-bit)
		\item RAM: 128 MB; 64 MB en Windows XP (32-bit)
		\item Espacio en disco: 124 MB
	\end{itemize}
	\item OS X
	\begin{itemize}
		\item Procesador Intel
		\item Mac OS X 10.7.3 (Lion) o posterior
		\item Privilegios de administrador para la instalación
	\end{itemize}
	\item Linux
	\begin{itemize}
		\item Oracle Linux 5.5+ / Red Hat Enterprise Linux 5.5+ / Ubuntu Linux 10.04+ / Suse Linux Enterprise Server 10 SP2+
	\end{itemize}
\end{itemize}

Los requerimientos del sistema vienen dados por la instalación de Java\footnote{\url{http://java.com/en/download/help/sysreq.xml}}.

\subsection{Instalación y ejecución}
\subsubsection{Instalación}
La aplicación se compone de un único archivo \emph{jar} ejecutable.
Por lo tanto, el método de instalación se reduce a copiar y pegar el fichero a la localización deseada.

\subsubsection{Ejecución}
Existen dos métodos posibles de ejecución.
\begin{itemize}
	\item Directamente haciendo \emph{doble-click} sobre el fichero \emph{jar}.
	Asumiendo que el \emph{path} de Java se encuentre correctamente configurado, el programa comenzará a ejecutarse inmediatamente.
	\item Mediante línea de comandos, navegando hasta la carpeta que contiene el fichero \emph{jar} y ejecutando el comando
	\begin{verbatim}
	java -jar PLQuiz-[versión].jar
	\end{verbatim}
\end{itemize}

\subsection{Resolución de problemas}
La aplicación permite la evaluación de conocimientos sobre dos tipos de algoritmos, mediante la composición de cuestionarios.
Cada uno de estos tipos dispone de tres sub-tipos adicionales, que se centran en aspectos distintos de un mismo problema general.
Los tipos y sub-tipos disponibles son los siguientes:
\begin{itemize}
	\item Evaluación de conocimientos de aplicación del algoritmo de Aho-Sethi-Ullman
	\begin{itemize}
		\item Evaluación de la construcción de árbol sintáctico.
		El objetivo es construir el árbol sintáctico correspondiente a la expresión regular dada, o elegir el correcto de entre los proporcionados.
		\item Evaluación de etiquetado del árbol.
		Tienen como objetivo el etiquetado de cada nodo del árbol sintáctico correspondiente a una expresión regular con sus correspondientes conjuntos \emph{primerapos} y \emph{últimapos}, indicando además si son o no anulables.
		\item Preguntas sobre construcción de tablas.
		El sub-tipo de cuestión más completo dentro de esta clase, asume la ejecución de los otros dos sub-tipos como paso previo a su resolución.
		Tiene como objetivo rellenar las tablas \emph{siguientepos} y de transición.
	\end{itemize}
	\item Preguntas sobre la aplicación del algoritmo de McNaughton-Yamada-Thompson
	\begin{itemize}
		\item Evaluación de la construcción del autómata.
		Tiene como objetivo construir el autómata finito correspondiente con la expresión regular dada, o elegir el correcto de entre los dados.
		\item Preguntas sobre la resolución de la expresión.
		El objetivo de este sub-tipo de cuestiones es la obtención de una tabla de transición a partir de una expresión regular dada.
		Asume la construcción del autómata como paso intermedio, combinando los otros dos sub-tipos de problema.
		\item Preguntas sobre resolución del autómata.
		Consiste en obtener la tabla de transición de un autómata finito dado.
	\end{itemize}
\end{itemize}

\subsubsection{Añadir una pregunta}
El añadido de problemas se realiza mediante los controles situados en la parte inferior izquierda de la ventana.
La lista desplegable permite seleccionar el tipo de problema deseado.

Una vez escogido el tipo, el botón marcado con `+' añadirá un problema de ese tipo al final de la lista de problemas actuales, situada en la parte izquierda de la ventana.

\imagen{manual_anadir}{Controles de añadido de problemas}

Existe la posibilidad de añadir bloques de múltiples problemas de cualquiera de los tipos al mismo tiempo, generando automáticamente sus contenidos.
El método para realizarlo se detalla en la sección \emph{generación de problemas}.

\subsubsection{Formato de la expresión regular}
La expresión regular debe concordar con un formato específico para ser correctamente reconocida por el parser.
\begin{itemize}
	\item Todo paréntesis de apertura `(' debe ir acompañado por un paréntesis de cierre `)'.
	\item Los símbolos utilizados deben ser una letra minúscula (a-z) o el símbolo `\$'.
	\item El carácter que representa la cadena vacía ($ \epsilon $) se representa con la letra `E' o con el carácter unicode \emph{\\u03B5}.
	\item La operación de cierre se representa con el símbolo `*', y va siempre precedida de un símbolo o de $ \epsilon $.
	\item La operación de concatenación se representa con el símbolo `.' o con el carácter unicode \emph{\\u2027}.
	\item La operación de unión se representa con el símbolo `$ | $'.
	\item Se admiten espacios en blanco en el interior de la expresión.
\end{itemize}

\subsubsection{Eliminar una pregunta}
La eliminación de un problema concreto se realiza mediante el botón situado a la izquierda de la caja de texto que contiene la expresión regular, identificado con el símbolo `-'.

\imagen{manual_eliminar}{Control de eliminación de preguntas}

La otra opción de eliminación de preguntas permite eliminar el cuestionario completo y empezar de cero.
Resulta accesible desde el menú `Archivo', tomando la opción de `Documento en blanco'.

\imagen{manual_nuevo}{Control de generación de nuevo documento en blanco}

Es importante remarcar que la eliminación de las preguntas es \emph{definitiva}.

\subsubsection{Tipos y sub-tipos de problema}

\subsection{Generación de preguntas}

\subsubsection{Características de la preguntas}
A la hora de generar una pregunta tenemos la opción de definir ciertas características del mismo, incluyendo:
\begin{itemize}
	\item Si la expresión regular asociada incorporará $ \epsilon $.
	\item El número de símbolos contenidos en la expresión regular, empezando por la `a' y sin incluir $ \epsilon $.
	Por ejemplo, la expresión $ ((a|b)*c)|\epsilon $ contendría tres símbolos.
	\item El número de estados en la tabla de transición del problema resuelto.
\end{itemize}
Las características se definen en la parte inferior del recuadro del problema.

\imagen{manual_caracteristicas}{Características seleccionadas y pregunta generada}

\subsubsection{Generación de expresiones}
Una vez definidas las características de la pregunta deseada, pulsar el botón `Generar' comenzará la búsqueda de una expresión regular que produzca dicha pregunta.
Mientras la aplicación realiza la búsqueda, se mostrará una barra de progreso en la parte inferior del recuadro de la pregunta, y el botón de `Generar' se transformará en el botón de `Cancelar'.
Dicho botón puede usarse durante la búsqueda para cancelar la operación.

\imagen{manual_generacion}{Generación de pregunta en curso}

Una vez generado la pregunta, la solución se muestra automáticamente de acuerdo con el sub-tipo de pregunta seleccionado.

La generación de preguntas se realiza en un proceso aislado, por lo que es posible interactuar con la interfaz durante el proceso, o generar múltiples preguntas al mismo tiempo.

\subsubsection{Generación de bloques de preguntas}
La aplicación no está limitada a la generación de preguntas individuales, sino que puede generar un número arbitrario de preguntas de cualquier tipo.
Para evitar la generación de un gran número de preguntas similares, la aplicación permite introducir un rango de variación, dentro de los cuales se encontrarán las características de las preguntas generados.

El acceso a la herramienta de generación de preguntas en bloque se encuentra en el menú `Archivo', seleccionando la opción `Generar bloque de problemas'.

\imagen{manual_menu_bloque}{Opción de menú para generación de preguntas en bloque}

Dentro de la interfaz de generación de preguntas podemos especificar el número de preguntas de cada tipo principal que queremos generar.
Cada tipo de preguntas permitirá seleccionar un sub-tipo, al cuál pertenecerán todas.

La selección de características deseadas tiene los mismos efectos que en la interfaz principal, con la diferencia de que permite especificar una variación.
Las preguntas generados en este bloque se ajustarán al rango definido entre el valor dado a la característica menos la variación, y el valor más la variación.

\imagen{manual_bloque}{Interfaz de generación de preguntas en bloque}

Al igual que la generación de preguntas normal, la generación en bloque muestra una barra de progreso y permite la cancelación.

Una limitación de este método es que todos las preguntas de un tipo dado que se generen pertenecerán al mismo sub-tipo.

\subsection{Exportación de cuestionarios}
La exportación del cuestionario actual se realiza mediante el menú `Exportar' de la barra de herramientas, seleccionando el formato deseado.
El cuestionario exportado no mostrará siempre el mismo aspecto que en la vista previa, y puede variar según el formato, de manera que la presentación sea la más adecuada en cada caso.

\imagen{manual_exportar}{Menú exportar, con todas las opciones de exportado de cuestionarios}

\subsubsection{Vista previa}
A diferencia del resto de formatos, la traducción de cuestionarios al formato de vista previa no está disponible al usuario, y se utiliza únicamente para la vista que aparece en el lado derecho de la interfaz gráfica.
El formato de salida es \emph{HTML} con estilos \emph{CSS} (aunque limitados), aprovechando la capacidad de los componentes \emph{Swing} para mostrarlo.

Dado que el contenido de un mismo cuestionario puede variar para los distintos formatos de salida, la vista previa mostrará siempre aquel que muestre más información.
Por ejemplo, si un cuestionario puede exportarse mostrando varias alternativas posibles además de la solución real, y con únicamente la solución, la vista previa mostrará las alternativas.

\subsubsection{Moodle \emph{XML}}
Esta traducción genera un cuestionario \emph{cloze} en formato \emph{XML} compatible con la plataforma de aprendizaje virtual Moodle.
El objetivo de un cuestionario \emph{cloze} es la elección de la solución correcta de entre varias posibles para cada pregunta.

Las imágenes de los cuestionarios Moodle no se almacenan de manera separada, sino que se insertan directamente en el documento \emph{XML}.
Para ello codificamos la imagen en forma de cadena de caracteres de 64 bits, y la incrustamos con etiquetas de archivo dentro del documento.

El importado del cuestionario \emph{cloze} se realiza en el menú `Banco de preguntas' de Moodle, sección `importar', y eligiendo el formato `Formato Moodle XML'.

\imagen{manual_importar_moodle}{Menú de importación de cuestionarios Moodle}

Una vez completado el proceso de importado, podemos ver las cuestiones individuales que se han añadido al banco de preguntas.
La aplicación marca cada cuestión generada con una identificación, indicando el tipo y sub-tipo de cada una.
Desde este menú podemos revisar las cuestiones, accediendo a una vista previa de las mismas, y eliminarlas del bando de preguntas.

\imagen{manual_moodle_banco}{Vista de cuestiones añadidas al banco de preguntas de Moodle}

\subsubsection{\LaTeX{}}
Los documentos formato \LaTeX{} pretenden ser examenes o cuestionarios imprimibles, por lo cual su formato tiende a ser diferente al de los cuestionarios virtuales.
La mayor diferencia es que los cuestionarios \emph{XML} se centran siempre en ofrecer varias opciones de respuesta, mientras que un cuestionario escrito puede exigir directamente la resolución del problema.

Por defecto el formato \LaTeX{} genera imágenes en formato \emph{JPG}, guardando el documento y las imágenes de manera separada en el directorio indicado.
El documento incrustará las imágenes directamente al ser compilado.
Los nombres dados a las imágenes son automáticos y únicos para cada imagen.

\subsubsection{\LaTeX{} con imágenes \emph{Graphviz}}
Este modo de exportado es idéntico al modo \LaTeX{} por defecto, salvo en el modo en que exporta las imágenes.
En lugar de generar la imagen directamente, lo que se exportan son programas en formato \emph{dot} para cada imagen, que pueden compilarse a \emph{pdf} o a otros formatos utilizando la herramienta \emph{Graphviz}.

Una vez generadas las imágenes, y sin importar el formato en que \emph{Graphviz} las emita, la compilación del documento \LaTeX{} es capaz de incluirlas directamente.

\section{Manual del programador}

\subsection{Requerimientos}
La preparación del proyecto para su desarrollo y empaquetado requiere obligatoriamente las siguientes herramientas:
\begin{itemize}
	\item Entorno de desarrollo Java 8 (Java JDK 8)\footnote{\url{http://www.oracle.com/technetwork/java/javase/downloads/jdk8-downloads-2133151.html}}
	\item Apache Maven\footnote{\url{ttp://maven.apache.org/}} versión 3 o posterior
\end{itemize}

Opcionalmente puede trabajarse con Eclipse como \emph{IDE}, con el complemento de \emph{JavaCC}\footnote{\url{http://eclipse-javacc.sourceforge.net/}}.
El complemento de Maven es opcional.
Java 8 exige Eclipse versión Luna o posterior.

\subsection{Preparación del proyecto}
\subsubsection{Descarga desde control de versiones}
El proyecto puede obtenerse directamente o a traves de su repositorio de Github\footnote{\url{https://github.com/RobertoIA/PLQuiz}}. Puede descargarse como un fichero comprimido mediante la opción \emph{Download ZIP}, utilizando uno de los clientes propietarios de Github con la opción \emph{Clone in Desktop}, o mediante la línea de comandos de git utilizando la URL de clonado\footnote{\url{https://github.com/RobertoIA/PLQuiz.git}} (ver figura~\ref{fig:manual_github}).

\imagenflotante{manual_github}{Opciones de descarga desde control de versiones}

\subsubsection{Resolución de dependencias}
La resolución de dependencias se realiza directamente mediante Maven.
Si no va a utilizarse Eclipse como IDE, puede usarse el siguiente comando en el directorio que contenga el fichero \emph{pom.xml}:
\begin{verbatim}
	mvn dependency:resolve
\end{verbatim}
En caso de utilizarse Eclipse, este comando resulta redundante, y el procedimiento a seguir se detalla en la versión siguiente.

\subsubsection{Trabajando con Eclipse}
En caso de utilizarse Eclipse como entorno de desarrollo, Maven nos permite realizar la resolución de dependencias y la conversión del proyecto a proyecto de Eclipse en un solo paso.
Esto se realiza con el complemento de Eclipse para Maven (y no viceversa), que Maven se encargará de descargar automáticamente si no hemos utilizado antes.

El siguiente comando debe ejecutarse en el directorio que contenga el fichero \emph{pom.xml}:
\begin{verbatim}
	mvn eclipse:eclipse
\end{verbatim}
Una vez ha terminado, podemos importar el proyecto desde eclipse con el menú `File', opción `Import'.
La opción a utilizar es `Existing Projects into Workspace', a la que solo tenemos que indicarle donde se encuentra almacenado el proyecto.

\subsubsection{Compilación y empaquetado}
La preparación de la aplicación para su uso se realiza con Maven, que se encarga de controlar las versiones, empaquetar las bibliotecas necesarias y realizar los test.
Para realizar el empaquetado de manera normal, ejecutaremos el siguiente comando en el directorio raíz del proyecto, donde encontremos el fichero \emph{pom.xml}.
\begin{verbatim}
	mvn package
\end{verbatim}
En caso de querer saltarnos la realización de test podemos utilizar el siguiente comando, en el mismo directorio que el anterior.
\begin{verbatim}
	mvn package -Dmaven.test.skip=true
\end{verbatim}
El fichero \emph{.jar} generado se encontrará en la carpeta \emph{target}, situada en la raíz del proyecto.

\subsection{Entorno Moodle}
La utilización de los cuestionarios exportados en formato Moodle XML requiere de la instalación de un entorno Moodle.
Los requerimientos del entorno son los siguientes:
\begin{itemize}
	\item Sistema operativo OS X, Windows o Linux
	\item 160MB de espacio libre en disco
	\item Procesador de 1GHZ
\end{itemize}

El proceso de instalación\footnote{\url{http://docs.moodle.org/27/en/Installing\_Moodle}} es relativamente complejo, incluyendo la instalación y configuración de Apache, MySQL y PHP.
Afortunadamente, Moodle ofrece instaladores completos para Windows\footnote{\url{http://download.moodle.org/windows/}} y OS X\footnote{\url{http://download.moodle.org/macosx/}} que se encargan de la instalación y configuración del entorno completo.
\apendice{Planificación por \emph{sprint}} \label{sec:Cappendix}

Esta sección detalla el trabajo realizado en cada \emph{sprint} del proyecto. Cada \emph{sprint} dura dos semanas, y corresponde con una serie de objetivos planteados y unos resultados conseguidos.
Al finalizar se realiza la entrega de una versión funcional de la aplicación.

Se incluye adicionalmente un \emph{sprint} inicial que transcurre entre la asignación del proyecto y el comienzo oficial.

\section{Prototipo --- hasta el 28 de febrero}

\subsection{Objetivos}
Desarrollo de un prototipo básico de la aplicación, como base a partir de la cual realizar las decisiones sobre el diseño de interfaz y la funcionalidad de la aplicación final.

El prototipo debe incluir una interfaz gráfica simple, y debe permitir resolver al menos uno de los algoritmos planteados. 
El código se desarrolla con la única intención de demostrar funcionalidad, y será desechado casi en su totalidad al final de la iteración.

\subsection{Resultados}
Al final de esta iteración contamos con una aplicación de interfaz gráfica que incluye un primer esbozo de los elementos que incluirá la aplicación final.
Algunos de estos elementos no son funcionales, pero se incluyen para referenciar funcionalidad futura.

El prototipo es capaz de procesar expresiones regulares a partir de cadenas de caracteres, extraer un árbol sintáctico de la misma y utilizarlo para obtener las tablas de siguiente posición y transiciones y demás datos correspondientes con el algoritmo de Aho-Sethi-Ullman.

Conservaremos (con modificaciones) el procesador de texto y la estructura de datos utilizada para el almacenamiento de expresiones regulares.

\section{Versión 0.1 --- 28 de febrero a 14 de marzo}
\subsection{Objetivos}
Diseñar y construir una aplicación como base para el proyecto, sin reutilizar código del prototipo pero con las mismas funcionalidades, e incluyendo documentación y pruebas.
Se realizarán modificaciones en la arquitectura general de la aplicación para facilitar la generación de ejercicios.

Se propone la idea de implementar un generador de ejercicios aleatorios basado en un algoritmo genético.

\subsection{Desarrollo}
Al ser la primera iteración realizada, carecemos de una estimación del progreso.
Realizando la revisión de la iteración destacamos que es la segunda iteración con más puntos completados, debido al gran número de tareas separadas que acompañan al comienzo del proyecto.

\imagen{iteracion_1}{Desarrollo de la primera iteración}

\subsection{Resultados}
Contamos con una implementación completa y estable del procesador de expresiones regulares y del algoritmo de Aho-Sethi-Ullman.
La interfaz es funcional y está preparada para ser ampliada con elementos posteriores.

Se añade funcionalidad de generación de cuestiones de tipo Aho-Sethi-Ullman implementando un algoritmo de búsqueda aleatoria.
La generación se realiza en función a una serie de parámetros dados por el usuario, que definiran ciertas características de la expresión regular y de los resultados de la cuestión.

Añadimos la base para el exportado de documentos en dos formatos diferentes: HTML y XML compatible con la plataforma Moodle.

\section{Versión 0.2 --- 14 de marzo a 28 de marzo}
\subsection{Objetivos}
Implementación del algoritmo de Thompson para la resolución de problemas de expresiones regulares.
Esto implica una serie de tareas separadas, como la construcción de autómatas finitos a partir de expresiones regulares, el añadido del tipo de problemas a la interfáz gráfica o su compatibilidad con el sistema de documentos.

Se pretende también reestructurar la construcción y el almacenamiento de documentos, tomando el trabajo de la iteración anterior como un primer prototipo.
También se toma con objetivo el añadido de documentos en formato \LaTeX{}.

\subsection{Desarrollo}
Sobrepasamos ampliamente la estimación de tareas a completar, siendo esta la iteración con más puntos completados de todo el desarrollo.
Se aprovecha el tiempo sobrante para refinar las tareas completadas, avanzando el proyecto todo lo posible.

\imagen{iteracion_2}{Desarrollo de la segunda iteración}

\subsection{Resultados}
El algoritmo de Thompson requeriría capacidades de dibujado de grafos para ser implementado correctamente.
Se implementa un tipo de cuestión relacionado, el de construcción de subconjuntos, a partir del cual implementaremos el de Thompson más adelante.

Completa la generación de documentos latex y en formato Moodle XML.
Se incluye la funcionalidad encargada de generar múltiples opciones de respuesta para las cuestiones que vayamos a exportar a Moodle.

A pesar de que no estaba planeado, el tiempo sobrante de esta iteración lo dedicamos a preparar un algoritmo de búsqueda aleatoria que permita la generación de cuestiones de construcción de subconjuntos, similar al desarrollado en la iteración anterior para el algoritmo de Aho-Sethi-Ullman.
Se comprueba que este algoritmo no resulta eficiente para este tipo de cuestiones, y se plantea la implementación de un algoritmo genético en una iteración posterior.

\section{Versión 0.3 --- 28 de marzo a 11 de abril}
\subsection{Objetivos}
Se plantean mejoras sobre el contenido de los ejercicios de Aho-Sethi-Ullman y construcción de subconjuntos.
Asimismo, se plantea el uso de algoritmos genéticos para la generación de cuestiones de construcción de subconjuntos.

\subsection{Desarrollo}
Esta iteración completa muy pocas tareas, siendo en general la menos productiva del desarrollo.
Se trabaja en completar ciertas partes de la aplicación desarrolladas en la iteración anterior, pero las tareas más complejas se retrasan hasta la iteración siguiente.

\imagen{iteracion_3}{Desarrollo de la tercera iteración}

\subsection{Resultados}
Se implementan los cambios solicitados para los ejercicios, generando las opciones de selección para el formato XML siguiendo el método de una opción correcta, dos opciones similares y una opción obviamente incorrecta.
Añadidas las bases para la implementación de algoritmos genéticos, como la posibilidad de realizar operaciones de mutación sobre las expresiones regulares.
Se prepara la implementación de generación concurrente de cuestiones, sin completar la funcionalidad.

\section{Versión 0.4 --- 11 de abril a 25 de abril}
\subsection{Objetivos}
Se plantea en primer lugar completar el sistema de generación de cuestiones, con los siguientes objetivos:
\begin{itemize}
	\item Permitir que la generación se complete de manera concurrente, sin bloquear la interfaz gráfica
	\item Crear un algoritmo capaz de generar ejercicios de construcción de subconjuntos, posiblemente un algoritmo genético
	\item Desarrollar un sistema que permita la generación de bloques de ejercicios
\end{itemize}

\subsection{Desarrollo}
Esta iteración completa una considerable cantidad de puntos, aunque el total de puntos pendientes apenas cambia debido a las nuevas tareas añadidas.
Se completan objetivos adicionales, como la funcionalidad que permite visualizar autómatas finitos en forma de imagen en la interfaz gráfica.
Estos añadidos producen nuevos objetivos, principalmente mejoras, correcciones y ampliaciones a la nueva funcionalidad.

\imagen{iteracion_4}{Desarrollo de la cuarta iteración}

\subsection{Resultados}
Se completan todos los objetivos de generación de ejercicios, dando ese subsistema por terminado al margen de retoques o modificaciones menores.

Como añadido se prepara el sistema que permite visualizar imágenes en la interfaz gráfica, y se desarrolla la funcionalidad de generar imágenes que representen los autómatas generados por la resolución de ejercicios de aplicación del algoritmo de McNaughton-Yamada-Thompson.

\section{Versión 0.5 --- 25 de abril a 9 de mayo}
\subsection{Objetivos}
Como objetivos de esta iteración marcamos en general la mejora de las características ya existentes, y el completado de la funcionalidad que está implementada solo parcialmente.
Se sugiere, por ejemplo, la inserción de imágenes en documentos XML mediante la codificación de la imagen a cadena de carácteres de 64 bits.

\subsection{Desarrollo}
Podemos ver como se completan los objetivos de esta iteración, con una subida de tareas pendientes únicamente cuando se añaden nuevas tareas al finalizar el periodo de desarrollo.

\imagen{iteracion_5}{Desarrollo de la quinta iteración}

\subsection{Resultados}
Este \emph{sprint} completa varios ejercicios adicionales, incluyendo mejoras en la visualización de imágenes de cara a la interfaz, la insercion de imágenes en documentos XML utilizando codificación de 64 bits, y en documentos \LaTeX{} guardándolas como ficheros.
También corrige problemas con las imágenes generadas, que hacían que no fueran completamente correctas.

\section{Versión 0.6 --- 9 de mayo a 23 de mayo}
\subsection{Objetivos}
Se plantea para esta iteración la ampliación de tipos de cuestiones disponibles, en concreto ejercicios de construcción de árboles sintácticos a partir de expresiones regulares, relacionado con la resolución del algoritmo de Aho-Sethi-Ullman.
Esto incluye también la generación de imágenes de árboles.

\subsection{Desarrollo}
La cantidad de puntos completados en esta iteración es irregular, ya que las tareas son amplias y no se dan por completadas hasta cerca del final, coincidiendo con el añadido de las nuevas tareas.

\imagen{iteracion_6}{Desarrollo de la sexta iteración}

\subsection{Resultados}
Completa el nuevo tipo de ejercicios.
También completa la integración de los nuevos tipos de ejercicios con el generador de bloques de cuestiones, permitiendo generar cuestiones del tipo y sub-tipo deseados.

\section{Versión 0.7 --- 23 de mayo a 6 de junio}
\subsection{Objetivos}
Se marca como objetivo la creación de un sistema de plantillas, que permita la generación de cuestionarios de acuerdo a un formato personalizable, e independiente para cada tipo de cuestión.
También se trabajará en ampliar la variedad de cuestiones disponibles.

\subsection{Desarrollo}
Se completan las tareas asignadas satisfactoriamente, así como las tareas añadidas durante el propio \emph{sprint}.
Se intenta completar las partes principales del desarrollo, de manera que puedan dedicarse las iteraciones posteriores al desarrollo de la memoria.

\imagen{iteracion_7}{Desarrollo de la séptima iteración}

\subsection{Resultados}
Se completa un nuevo tipo de cuestión, y se da por finalizado el sistema de plantillas.
Damos por terminado el desarrollo del proyecto, salvo detalles y retoques.

\section{Versión 0.8 --- 6 de junio a 20 de junio}
\subsection{Objetivos}
En este \emph{sprint} se comienza a desarrollar la memoria del proyecto a tiempo completo.
Como actividad adicional se plantea la posibilidad de generar imágenes en un formato extra, compatible con la herramienta Graphviz.

\subsection{Desarrollo}
Se completan una cantidad satisfactoria de objetivos.
Podemos ver que quedan pendientes partes de la memoria para las iteraciones posteriores, en forma de tareas adicionales.

\imagen{iteracion_8}{Desarrollo de la octava iteración}

\subsection{Resultados}
Se completa la parte inicial de la memoria, incluyendo la descripción de herramientas utilizadas en el desarrollo y el apéndice en el que se analiza la generación de expresiones regulares.

\section{Versión 0.9 --- 20 de junio a 4 de julio}
\subsection{Objetivos}
El objetivo de esta iteración es avanzar en el desarrollo de la memoria del proyecto, aplicando a la aplicación los retoques que se vean necesarios pero sin añadir nueva funcionalidad.

\subsection{Desarrollo}
El número de puntos asignados a tareas y el propio número de tareas sufre variaciones, al seleccionarse que partes se incluyen en la memoria.

\imagen{iteracion_9}{Desarrollo de la novena iteración}

\subsection{Resultados}
Se avanza en el desarrollo de la memoria, incluyendo la parte de conceptos teóricos.
También se añade un nuevo tipo de cuestión a la aplicación, consistente en emparejar una expresión regular con el autómata finito que resulta de aplicarle el algoritmo de McNaughton-Yamada-Thompson.

\section{Versión 1.0 --- 4 de julio a 18 de julio}
\subsection{Objetivos}
La reunión de este último \emph{sprint} se adelanta, al coincidir de otra manera con la fecha de finalización del proyecto.
Los objetivos son la finalización de la memoria, y la preparación de la entrega final.

\subsection{Desarrollo}
El gráfico \emph{burn-down} de esta iteración termina antes de lo esperado, dado que el completar este informe entra dentro del propio plazo.
Este \emph{sprint} completa todas las tareas restantes.

\imagen{iteracion_10}{Desarrollo de la décima iteración}

\subsection{Resultados}
Se completa la versión final de la memoria, y se preparan todos los materiales para la entrega.
Con este \emph{sprint} se da por finalizado el proyecto.

\section{Resumen del proceso de desarrollo}
Observando la figura \ref{fig:puntos_desarrollo} vemos que el número de puntos aceptados en cada \emph{sprint} es altamente irregular.
Esto se debe al distinto número de tareas en cada sprint, dado que muchas tareas pequeñas completadas tienden a contar más que pocas tareas grandes.
Asimismo se ve afectado por el avance del curso, con una reducción del trabajo contribuido en las épocas de más actividad.

\imagen{puntos_desarrollo}{Puntos de tarea completados por iteración}

Podemos ver que las primeras iteraciones son las más fructiferas, dado que consiste en sentar las bases del proyecto.
Estas iteraciones incluyen una cantidad alta de pequeñas tareas muy diversas, pero relativamente simples de completar.

Las últimas iteraciones tienden a tener un número similar ---y más bajo--- de puntos.
Estos \emph{sprints} están dedicados a la redacción de la memoria, y por tanto constan de tareas amplias (secciones o capítulos completos), que tienen un alto valor en puntos pero que tienden a contribuir menos al total.

\bibliographystyle{plain}
\bibliography{PLQuiz}

\end{document}